\section{Discussion and lessons learned}
\label{sec:discussion}

The lifecycle audits show what a collection-level representation buys in
practice. A rename updated document state while leaving nine stale chunks
in the vector index. Re-indexing removed the duplicates and restored the
expected count. The three-source checks let operators verify the file,
registry and metadata views separately, while a search probe checks whether
the index reflects the change. The paper-ingestion audit reports separate
observations for a derived library and the five-library inventory
(Section~\ref{sec:e1}). The practical implication is to
test retrieval after a lifecycle operation as well as checking its write
response. This joins metadata governance to a property that matters at
query time \citep{park-metadata-quality-2009,paivarinta-ecm-2005}.

At the access layer, the full route cards gave the scripted selector a
usable collection-level view. The evidence packs show a second decision
after routing: how much of the selected material to read. Adding document
and vector chunk windows in A0h raised target-string containment over A0 by
15 questions in the frozen run and 12 in each same-question rerun. The
frozen run used an average of 309~kB of trace-recorded text per question,
compared with 31.4~kB for A0 and 24.3~MB for the mechanical full read. Thus
A0h traded a larger reading budget for more retained targets while using
about 1.3\% of the full scan's text volume. A0h also led the per-library
union PD by six questions in both the frozen and November~12 records, with a smaller paired difference
that remained unresolved. PD searches each of the ten libraries separately
before merging document candidates. The stronger result is the improvement over A0
when the retrieval configuration expands the evidence it reads.

Source visibility carries through to answering. On different substrates,
the rates of detected document-style markers diverged across the three
interfaces. Counts were nine of ten for each arm in the repeat using the same
indexed libraries, but the changed BQ06 wording and access configurations
prevent attribution of that convergence to the substrate alone. Marker checks are a
trace-format diagnostic. A separate semantic assessment would be needed to
determine which answer claims the retrieved passages support.

The historical ordering comparison identifies a distinct design choice.
Among retained candidates, fusion order produced 15 top-five hits across
sixteen queries, versus seven in verifier-score order. The separate score
archive contains only scores above the threshold; it does not establish the
effect of active filtering. This distinction parallels the practical
separation of relevance judgement and ranking in work on
LLM judgements and reranking \citep{zheng-llmjudge,sun-urara}.

Operational traces also show why a search needs to report its coverage.
Pagination recovered a failed catalogue path. An index read failure exposed
a different defect: the global channel could skip a library and still report
success. For the affected example, re-indexing recovered the missing hit.
Reporting partial coverage would help callers distinguish an incomplete search
from a completed search that returned no supporting evidence.
